\usepackage{etoolbox}
\usepackage[most]{tcolorbox}
\usepackage{graphicx}
\usepackage{xparse}
\tcbuselibrary{breakable}

% --- Disable floating figures globally ---
\usepackage{float}
\let\origfigure\figure
\let\endorigfigure\endfigure
\renewenvironment{figure}[1][H]{\origfigure[H]}{\endorigfigure}

% --- Colors ---
\usepackage{xcolor}
\definecolor{speculativeback}{RGB}{247,232,255}  % #f7e8ff
\definecolor{speculativeframe}{RGB}{139,43,226}  % #8b2be2

% --- Left-align all Pandoc tables ---
\usepackage{longtable}
\usepackage{etoolbox}
\makeatletter
\patchcmd\longtable{\par}{\par\raggedright}{}{}
\setlength{\LTleft}{0pt}
\setlength{\LTright}{0pt}
\makeatother

% Enable index generation
\usepackage{makeidx}
\makeindex


% Helper to show a blank title if empty
\newcommand{\blanktitle}{\mbox{}}

% =========================================================
% Concept Boxes
% =========================================================
\newtcolorbox{infobox}[1][Information]{%
  title=\ifstrempty{#1}{\blanktitle}{#1},
  colback=blue!3!white,
  colframe=blue!50!black,
  colbacktitle=blue!10!white,
  fonttitle=\bfseries,
  coltitle=black,
  enhanced, sharp corners, breakable,
  after=\par\vspace{6pt}
}

\newtcolorbox{proposedbox}[1][Proposed Concept]{%
  title=\ifstrempty{#1}{\blanktitle}{#1},
  colback=blue!5!white,
  colframe=blue!75!black,
  colbacktitle=blue!15!white,
  fonttitle=\bfseries,
  coltitle=black,
  enhanced, sharp corners, breakable,
  after=\par\vspace{6pt}
}

\newtcolorbox{establishedbox}[1][Established Concept]{%
  title=\ifstrempty{#1}{\blanktitle}{#1},
  colback=green!5!white,
  colframe=green!75!black,
  colbacktitle=green!15!white,
  fonttitle=\bfseries,
  coltitle=black,
  enhanced, sharp corners, breakable,
  after=\par\vspace{6pt}
}

\newtcolorbox{speculativebox}[1][Speculative Concept]{%
  title=\ifstrempty{#1}{\blanktitle}{#1},
  colback=speculativeback,
  colframe=speculativeframe,
  colbacktitle=speculativeback!80!white, % lighter variant for title bar
  fonttitle=\bfseries,
  coltitle=black,
  enhanced, sharp corners, breakable,
  after=\par\vspace{6pt}
}


\newtcolorbox{cautionbox}[1][Caution]{%
  title=\ifstrempty{#1}{\blanktitle}{#1},
  colback=orange!5!white,
  colframe=orange!75!black,
  colbacktitle=orange!15!white,
  fonttitle=\bfseries,
  coltitle=black,
  enhanced, sharp corners, breakable,
  after=\par\vspace{6pt}
}

\newtcolorbox{warningbox}[1][Warning]{%
  title=\ifstrempty{#1}{\blanktitle}{#1},
  colback=orange!5!white,
  colframe=orange!75!black,
  colbacktitle=orange!15!white,
  fonttitle=\bfseries,
  coltitle=black,
  enhanced, sharp corners, breakable,
  after=\par\vspace{6pt}
}

\newtcolorbox{dangerbox}[1][Danger]{%
  title=\ifstrempty{#1}{\blanktitle}{#1},
  colback=red!5!white,
  colframe=red!75!black,
  colbacktitle=red!15!white,
  fonttitle=\bfseries,
  coltitle=black,
  enhanced, sharp corners, breakable,
  after=\par\vspace{6pt}
}

\usepackage{xurl}

\usepackage{fancyhdr}
\setlength{\headheight}{15pt}

% --- Chapter-opening pages (plain style) ---
\fancypagestyle{plain}{%
  \fancyhf{}%
  \fancyhead[C]{\leftmark}      % Chapter title
  \cfoot{Basic Book Builder: A Pandoc template for building books and articles}   % Book title
  \rfoot{\thepage}              % Page number
  \renewcommand{\headrulewidth}{0pt}
  \renewcommand{\footrulewidth}{0pt}
}

% --- Normal running pages ---
\pagestyle{fancy}
\fancyhf{}

% Header
\fancyhead[C]{\leftmark}        % Chapter title (centered)

% Footer
\cfoot{Basic Book Builder: A Pandoc template for building books and articles}     % Book title
\rfoot{\thepage}                % Page number

% Line rules
\renewcommand{\headrulewidth}{0pt}   % rule under header
\renewcommand{\footrulewidth}{0pt}     % no rule above footer
